\section{Conclusion}
\label{sec:conclusion}

Reliable prognostics must remain coherent when the sensors available to an asset change. In a controlled synthetic evaluation with exact oracle posteriors, explicitly conditioning the predictor on the observed sensor set reduced cross-view posterior inconsistency by 38.6\% (90\% CI: 34.8--42.4\%) when the alternative views remained task-sufficient. The effect persisted after a four-fold reduction in Monte Carlo sampling, suggesting that the gain is tied to the representation of the observation process rather than to a large sampling budget.

The same study identifies the boundary of that gain. Under permanent removal of functionally critical sensors, PROP-DIRECT did not retain its advantage, while exploratory oracle-model correlations showed a weaker distinction between critical and random deletions than the standard baseline. These comparisons are descriptive because the cases share assets and query structure, but they point to a concrete architectural issue: recording which sensors are present is easier than representing what those sensors mean for the degradation process. The result is a useful design constraint for future importance-aware conditioning, rather than a reason to discard observation-process conditioning altogether.

On C-MAPSS FD001, the shared-model evaluation reduces distance to the full view by 18.5--28.9\%, but increases the distance between the two equally sized reduced views by 36.6\%. Point MAE remains close to the baseline. This conditional pattern extends the geometric result to a public simulated benchmark while sharpening the importance-awareness boundary; evaluation with naturally occurring sensor loss remains the next test.

Taken together, the findings support a simple principle for probabilistic prognostics: observation-process conditioning is valuable when available views preserve the information needed by the task, and it needs an explicit account of functional importance when they do not. Hierarchical sensor representations, state-dependent importance weights, and physics-informed priors offer natural ways to pursue that account while retaining the geometric objective introduced here.
